\documentclass[12pt,a4paper]{article}
\usepackage[utf8]{vietnam}
\usepackage{amsmath,amssymb}
\usepackage{tabularx}
\usepackage{xcolor}
\usepackage[top=0.6cm, bottom=0.6cm, left=0.6cm, right=0.6cm]{geometry}
\usepackage[dethi]{ex_test}

\setlength{\parskip}{0pt}

% Dòng kẻ chấm mờ chuẩn viết tay
\newcommand{\dongke}[1]{%
  \begingroup%
  \setlength{\parskip}{0pt}%
  \baselineskip=18.5pt%
  \vspace{1.0mm}%
  \foreach \i in {1,...,#1}{%
    \noindent{\color{gray!60}\dotfill}\par%
  }%
  \vspace{0.6mm}%
  \endgroup%
}

% Định nghĩa ô bài tập tự luận trong bảng
\newcommand{\oBaiTap}[3][6]{%
  \vspace{1.0mm}%
  \noindent\textbf{#2} #3\par%
  \dongke{#1}%
}

\begin{document}
\pagestyle{empty}
% =============================================================================
% ========================= MA DE 101 - MAT TRUOC (TRANG 1/2) =============
% =============================================================================
\noindent
\begin{tabularx}{\linewidth}{|X|c|}
\hline
\textbf{PHIẾU RÈN LUYỆN CHUYÊN ĐỀ TOÁN 6 - CÂU 11 \& CÂU 12} & \textbf{Mã đề: 101 - Trang 1/2} \\
\hline
\end{tabularx}

\vspace{1.0mm}

\noindent
\begin{tabularx}{\linewidth}{|p{3.0cm}|p{3.0cm}|X|}
\hline
\centering\textbf{Điểm Câu 11} & \centering\textbf{Điểm Câu 12} & \centering\textbf{Lời nhận xét của thầy cô} \tabularnewline
\rule{0pt}{1.8cm} & & \\
\hline
\end{tabularx}

\vspace{1.5mm}

\noindent\textbf{Bài 1 (2,0 điểm). Tính một cách hợp lý (nếu có thể):}
\vspace{0.5mm}

\noindent
\begin{tabularx}{\linewidth}{|X|X|}
\hline
\oBaiTap[6]{a)}{145 + 328 + 255 + 172} & \oBaiTap[6]{b)}{43 \cdot 78 + 43 \cdot 22} \\
\hline
\oBaiTap[6]{c)}{35 \cdot 118 + 35 \cdot 82 - 35 \cdot 100} & \oBaiTap[6]{d)}{(-25) + 84 + (-75) + 16} \\
\hline
\end{tabularx}

\vspace{1.5mm}

\noindent\textbf{Bài 2 (3,0 điểm). Thực hiện phép tính:}
\vspace{0.5mm}

\noindent
\begin{tabularx}{\linewidth}{|X|X|}
\hline
\oBaiTap[5]{a)}{3 \cdot 2^3 - 36 : 3^2 + 2026^0} & \oBaiTap[5]{b)}{150 : (38 - 13) + 3^2 \cdot 4} \\
\hline
\oBaiTap[5]{c)}{180 : [45 - (11 - 6)^2]} & \oBaiTap[5]{d)}{200 : \{60 - [45 - (12 - 7)^2]\}} \\
\hline
\oBaiTap[5]{e)}{5^7 : 5^5 + 2^3 \cdot 2^2 - 11^{2026} : 11^{2025}} & \oBaiTap[5]{f)}{120 - [30 - (4 - 1)^2] \cdot 5} \\
\hline
\end{tabularx}
\newpage
% =============================================================================
% ========================= MA DE 101 - MAT SAU (TRANG 2/2) ==============
% =============================================================================
\noindent
\begin{tabularx}{\linewidth}{|X|r|}
\hline
\textbf{TOÁN 6 - CHUYÊN ĐỀ CÂU 11 \& CÂU 12} & \textbf{Mã đề: 101 - Trang 2/2} \\
\hline
\end{tabularx}

\vspace{1.5mm}

\noindent\textbf{Bài 3 (2,0 điểm). Tìm số tự nhiên $x$, biết:}
\vspace{0.5mm}

\noindent
\begin{tabularx}{\linewidth}{|X|X|}
\hline
\oBaiTap[5]{a)}{3x - 17 = 28} & \oBaiTap[5]{b)}{5(x - 4) = 45} \\
\hline
\oBaiTap[5]{c)}{4(x + 5) - 18 = 42} & \oBaiTap[5]{d)}{85 - 3(x + 2) = 40} \\
\hline
\end{tabularx}

\vspace{1.5mm}

\noindent\textbf{Bài 4 (1,5 điểm). Tìm số tự nhiên $x$ chứa lũy thừa:}
\vspace{0.5mm}

\noindent
\begin{tabularx}{\linewidth}{|X|X|X|}
\hline
\oBaiTap[6]{a)}{3 \cdot 2^x = 48} & \oBaiTap[6]{b)}{4 \cdot 2^{x+1} - 15 = 49} & \oBaiTap[6]{c)}{(x - 3)^2 = 25 \quad (x \in \mathbb{N})} \\
\hline
\end{tabularx}

\vspace{1.5mm}

\noindent\textbf{Bài 5 (1,5 điểm). Tìm số nguyên $x$ \& Bài toán chia hết:}
\vspace{0.5mm}

\noindent
\begin{tabularx}{\linewidth}{|X|}
\hline
\oBaiTap[3]{a)}{Liệt kê các số nguyên $x$ thỏa mãn: $-4 \le x < 3$.} \\
\hline
\oBaiTap[4]{b)}{Cho các số nguyên $x$ thỏa mãn $-5 < x \le 6$. Tính tổng tất cả các số nguyên $x$.} \\
\hline
\oBaiTap[5]{c)}{Tìm số nguyên $x$ biết $(x - 2)$ là ước của $6$.} \\
\hline
\end{tabularx}
\newpage

% =============================================================================
% ========================= MA DE 102 - MAT TRUOC (TRANG 1/2) =============
% =============================================================================
\noindent
\begin{tabularx}{\linewidth}{|X|c|}
\hline
\textbf{PHIẾU RÈN LUYỆN CHUYÊN ĐỀ TOÁN 6 - CÂU 11 \& CÂU 12} & \textbf{Mã đề: 102 - Trang 1/2} \\
\hline
\end{tabularx}

\vspace{1.0mm}

\noindent
\begin{tabularx}{\linewidth}{|p{3.0cm}|p{3.0cm}|X|}
\hline
\centering\textbf{Điểm Câu 11} & \centering\textbf{Điểm Câu 12} & \centering\textbf{Lời nhận xét của thầy cô} \tabularnewline
\rule{0pt}{1.8cm} & & \\
\hline
\end{tabularx}

\vspace{1.5mm}

\noindent\textbf{Bài 1 (2,0 điểm). Tính một cách hợp lý (nếu có thể):}
\vspace{0.5mm}

\noindent
\begin{tabularx}{\linewidth}{|X|X|}
\hline
\oBaiTap[6]{a)}{136 + 247 + 264 + 153} & \oBaiTap[6]{b)}{57 \cdot 64 + 57 \cdot 36} \\
\hline
\oBaiTap[6]{c)}{48 \cdot 124 + 48 \cdot 76 - 48 \cdot 100} & \oBaiTap[6]{d)}{(-35) + 68 + (-65) + 32} \\
\hline
\end{tabularx}

\vspace{1.5mm}

\noindent\textbf{Bài 2 (3,0 điểm). Thực hiện phép tính:}
\vspace{0.5mm}

\noindent
\begin{tabularx}{\linewidth}{|X|X|}
\hline
\oBaiTap[5]{a)}{4 \cdot 2^3 - 45 : 3^2 + 2025^0} & \oBaiTap[5]{b)}{160 : (42 - 26) + 4^2 \cdot 3} \\
\hline
\oBaiTap[5]{c)}{240 : [55 - (13 - 8)^2]} & \oBaiTap[5]{d)}{240 : \{70 - [55 - (14 - 9)^2]\}} \\
\hline
\oBaiTap[5]{e)}{6^7 : 6^5 + 2^4 \cdot 2^1 - 13^{2025} : 13^{2024}} & \oBaiTap[5]{f)}{140 - [40 - (5 - 2)^2] \cdot 4} \\
\hline
\end{tabularx}
\newpage
% =============================================================================
% ========================= MA DE 102 - MAT SAU (TRANG 2/2) ==============
% =============================================================================
\noindent
\begin{tabularx}{\linewidth}{|X|r|}
\hline
\textbf{TOÁN 6 - CHUYÊN ĐỀ CÂU 11 \& CÂU 12} & \textbf{Mã đề: 102 - Trang 2/2} \\
\hline
\end{tabularx}

\vspace{1.5mm}

\noindent\textbf{Bài 3 (2,0 điểm). Tìm số tự nhiên $x$, biết:}
\vspace{0.5mm}

\noindent
\begin{tabularx}{\linewidth}{|X|X|}
\hline
\oBaiTap[5]{a)}{4x - 19 = 45} & \oBaiTap[5]{b)}{6(x - 5) = 48} \\
\hline
\oBaiTap[5]{c)}{5(x + 4) - 25 = 35} & \oBaiTap[5]{d)}{95 - 4(x + 1) = 35} \\
\hline
\end{tabularx}

\vspace{1.5mm}

\noindent\textbf{Bài 4 (1,5 điểm). Tìm số tự nhiên $x$ chứa lũy thừa:}
\vspace{0.5mm}

\noindent
\begin{tabularx}{\linewidth}{|X|X|X|}
\hline
\oBaiTap[6]{a)}{2 \cdot 3^x = 54} & \oBaiTap[6]{b)}{5 \cdot 2^{x+1} - 28 = 52} & \oBaiTap[6]{c)}{(x - 4)^2 = 36 \quad (x \in \mathbb{N})} \\
\hline
\end{tabularx}

\vspace{1.5mm}

\noindent\textbf{Bài 5 (1,5 điểm). Tìm số nguyên $x$ \& Bài toán chia hết:}
\vspace{0.5mm}

\noindent
\begin{tabularx}{\linewidth}{|X|}
\hline
\oBaiTap[3]{a)}{Liệt kê các số nguyên $x$ thỏa mãn: $-5 < x \le 2$.} \\
\hline
\oBaiTap[4]{b)}{Cho các số nguyên $x$ thỏa mãn $-6 \le x < 7$. Tính tổng tất cả các số nguyên $x$.} \\
\hline
\oBaiTap[5]{c)}{Tìm số nguyên $x$ biết $(x + 1)$ là ước của $8$.} \\
\hline
\end{tabularx}
\newpage

% =============================================================================
% ========================= MA DE 103 - MAT TRUOC (TRANG 1/2) =============
% =============================================================================
\noindent
\begin{tabularx}{\linewidth}{|X|c|}
\hline
\textbf{PHIẾU RÈN LUYỆN CHUYÊN ĐỀ TOÁN 6 - CÂU 11 \& CÂU 12} & \textbf{Mã đề: 103 - Trang 1/2} \\
\hline
\end{tabularx}

\vspace{1.0mm}

\noindent
\begin{tabularx}{\linewidth}{|p{3.0cm}|p{3.0cm}|X|}
\hline
\centering\textbf{Điểm Câu 11} & \centering\textbf{Điểm Câu 12} & \centering\textbf{Lời nhận xét của thầy cô} \tabularnewline
\rule{0pt}{1.8cm} & & \\
\hline
\end{tabularx}

\vspace{1.5mm}

\noindent\textbf{Bài 1 (2,0 điểm). Tính một cách hợp lý (nếu có thể):}
\vspace{0.5mm}

\noindent
\begin{tabularx}{\linewidth}{|X|X|}
\hline
\oBaiTap[6]{a)}{158 + 264 + 242 + 136} & \oBaiTap[6]{b)}{68 \cdot 83 + 68 \cdot 17} \\
\hline
\oBaiTap[6]{c)}{52 \cdot 135 + 52 \cdot 65 - 52 \cdot 100} & \oBaiTap[6]{d)}{(-45) + 73 + (-55) + 27} \\
\hline
\end{tabularx}

\vspace{1.5mm}

\noindent\textbf{Bài 2 (3,0 điểm). Thực hiện phép tính:}
\vspace{0.5mm}

\noindent
\begin{tabularx}{\linewidth}{|X|X|}
\hline
\oBaiTap[5]{a)}{5 \cdot 2^3 - 18 : 3^2 + 2024^0} & \oBaiTap[5]{b)}{210 : (51 - 30) + 5^2 \cdot 2} \\
\hline
\oBaiTap[5]{c)}{280 : [65 - (14 - 9)^2]} & \oBaiTap[5]{d)}{280 : \{80 - [65 - (13 - 8)^2]\}} \\
\hline
\oBaiTap[5]{e)}{4^7 : 4^5 + 3^2 \cdot 3^2 - 15^{2024} : 15^{2023}} & \oBaiTap[5]{f)}{160 - [50 - (6 - 3)^2] \cdot 3} \\
\hline
\end{tabularx}
\newpage
% =============================================================================
% ========================= MA DE 103 - MAT SAU (TRANG 2/2) ==============
% =============================================================================
\noindent
\begin{tabularx}{\linewidth}{|X|r|}
\hline
\textbf{TOÁN 6 - CHUYÊN ĐỀ CÂU 11 \& CÂU 12} & \textbf{Mã đề: 103 - Trang 2/2} \\
\hline
\end{tabularx}

\vspace{1.5mm}

\noindent\textbf{Bài 3 (2,0 điểm). Tìm số tự nhiên $x$, biết:}
\vspace{0.5mm}

\noindent
\begin{tabularx}{\linewidth}{|X|X|}
\hline
\oBaiTap[5]{a)}{5x - 23 = 52} & \oBaiTap[5]{b)}{7(x - 3) = 56} \\
\hline
\oBaiTap[5]{c)}{6(x + 3) - 22 = 38} & \oBaiTap[5]{d)}{105 - 5(x + 2) = 45} \\
\hline
\end{tabularx}

\vspace{1.5mm}

\noindent\textbf{Bài 4 (1,5 điểm). Tìm số tự nhiên $x$ chứa lũy thừa:}
\vspace{0.5mm}

\noindent
\begin{tabularx}{\linewidth}{|X|X|X|}
\hline
\oBaiTap[6]{a)}{5 \cdot 2^x = 80} & \oBaiTap[6]{b)}{3 \cdot 3^{x+1} - 25 = 56} & \oBaiTap[6]{c)}{(x - 2)^3 = 27 \quad (x \in \mathbb{N})} \\
\hline
\end{tabularx}

\vspace{1.5mm}

\noindent\textbf{Bài 5 (1,5 điểm). Tìm số nguyên $x$ \& Bài toán chia hết:}
\vspace{0.5mm}

\noindent
\begin{tabularx}{\linewidth}{|X|}
\hline
\oBaiTap[3]{a)}{Liệt kê các số nguyên $x$ thỏa mãn: $-3 \le x \le 4$.} \\
\hline
\oBaiTap[4]{b)}{Cho các số nguyên $x$ thỏa mãn $-7 < x \le 8$. Tính tổng tất cả các số nguyên $x$.} \\
\hline
\oBaiTap[5]{c)}{Tìm số nguyên $x$ biết $10 \ \vdots \ (x - 1)$.} \\
\hline
\end{tabularx}
\newpage

% =============================================================================
% ========================= MA DE 104 - MAT TRUOC (TRANG 1/2) =============
% =============================================================================
\noindent
\begin{tabularx}{\linewidth}{|X|c|}
\hline
\textbf{PHIẾU RÈN LUYỆN CHUYÊN ĐỀ TOÁN 6 - CÂU 11 \& CÂU 12} & \textbf{Mã đề: 104 - Trang 1/2} \\
\hline
\end{tabularx}

\vspace{1.0mm}

\noindent
\begin{tabularx}{\linewidth}{|p{3.0cm}|p{3.0cm}|X|}
\hline
\centering\textbf{Điểm Câu 11} & \centering\textbf{Điểm Câu 12} & \centering\textbf{Lời nhận xét của thầy cô} \tabularnewline
\rule{0pt}{1.8cm} & & \\
\hline
\end{tabularx}

\vspace{1.5mm}

\noindent\textbf{Bài 1 (2,0 điểm). Tính một cách hợp lý (nếu có thể):}
\vspace{0.5mm}

\noindent
\begin{tabularx}{\linewidth}{|X|X|}
\hline
\oBaiTap[6]{a)}{127 + 349 + 273 + 151} & \oBaiTap[6]{b)}{76 \cdot 75 + 76 \cdot 25} \\
\hline
\oBaiTap[6]{c)}{63 \cdot 142 + 63 \cdot 58 - 63 \cdot 100} & \oBaiTap[6]{d)}{(-15) + 59 + (-85) + 41} \\
\hline
\end{tabularx}

\vspace{1.5mm}

\noindent\textbf{Bài 2 (3,0 điểm). Thực hiện phép tính:}
\vspace{0.5mm}

\noindent
\begin{tabularx}{\linewidth}{|X|X|}
\hline
\oBaiTap[5]{a)}{6 \cdot 2^3 - 54 : 3^2 + 2027^0} & \oBaiTap[5]{b)}{240 : (46 - 22) + 3^2 \cdot 5} \\
\hline
\oBaiTap[5]{c)}{300 : [75 - (12 - 7)^2]} & \oBaiTap[5]{d)}{320 : \{90 - [75 - (15 - 10)^2]\}} \\
\hline
\oBaiTap[5]{e)}{7^7 : 7^5 + 2^2 \cdot 2^3 - 17^{2027} : 17^{2026}} & \oBaiTap[5]{f)}{180 - [60 - (7 - 4)^2] \cdot 3} \\
\hline
\end{tabularx}
\newpage
% =============================================================================
% ========================= MA DE 104 - MAT SAU (TRANG 2/2) ==============
% =============================================================================
\noindent
\begin{tabularx}{\linewidth}{|X|r|}
\hline
\textbf{TOÁN 6 - CHUYÊN ĐỀ CÂU 11 \& CÂU 12} & \textbf{Mã đề: 104 - Trang 2/2} \\
\hline
\end{tabularx}

\vspace{1.5mm}

\noindent\textbf{Bài 3 (2,0 điểm). Tìm số tự nhiên $x$, biết:}
\vspace{0.5mm}

\noindent
\begin{tabularx}{\linewidth}{|X|X|}
\hline
\oBaiTap[5]{a)}{6x - 26 = 58} & \oBaiTap[5]{b)}{8(x - 2) = 64} \\
\hline
\oBaiTap[5]{c)}{7(x + 2) - 26 = 30} & \oBaiTap[5]{d)}{115 - 6(x + 1) = 43} \\
\hline
\end{tabularx}

\vspace{1.5mm}

\noindent\textbf{Bài 4 (1,5 điểm). Tìm số tự nhiên $x$ chứa lũy thừa:}
\vspace{0.5mm}

\noindent
\begin{tabularx}{\linewidth}{|X|X|X|}
\hline
\oBaiTap[6]{a)}{4 \cdot 3^x = 108} & \oBaiTap[6]{b)}{2 \cdot 5^{x+1} - 18 = 32} & \oBaiTap[6]{c)}{(x - 1)^3 = 64 \quad (x \in \mathbb{N})} \\
\hline
\end{tabularx}

\vspace{1.5mm}

\noindent\textbf{Bài 5 (1,5 điểm). Tìm số nguyên $x$ \& Bài toán chia hết:}
\vspace{0.5mm}

\noindent
\begin{tabularx}{\linewidth}{|X|}
\hline
\oBaiTap[3]{a)}{Liệt kê các số nguyên $x$ thỏa mãn: $-6 < x < 2$.} \\
\hline
\oBaiTap[4]{b)}{Cho các số nguyên $x$ thỏa mãn $-5 \le x < 5$. Tính tổng tất cả các số nguyên $x$.} \\
\hline
\oBaiTap[5]{c)}{Tìm số nguyên $x$ biết $(x - 3)$ là ước của $9$.} \\
\hline
\end{tabularx}

\end{document}
